\section{Корпус дослідження}
\label{sec:corpus}

\subsection{Джерела та критерії відбору матеріалу}

Матеріалом дослідження є відеозаписи звернень, опубліковані на офіційному вебпорталі Офісу Президента (\texttt{president.gov.ua}) та відповідному YouTube-каналі. Вибір цього
джерела обумовлений кількома чинниками:
\begin{itemize}[nosep]
    \item усі записи виконано українською мовою, що відповідає напрямку
    перекладу, який досліджується;
    \item тематика звернень належить до сфери політичного та дипломатичного
    дискурсу, що характеризується складною лексикою, специфічною термінологією
    та високими вимогами до точності перекладу;
    \item записи мають високу якість аудіо (студійний запис, один мовець,
    відсутність стороннього шуму), що зменшує вплив технічних артефактів
    на результати розпізнавання мовлення.
\end{itemize}

Критерії відбору конкретних записів:
\begin{enumerate}[nosep]
    \item тривалість запису --- від 5 до 15 хвилин (для отримання фрагментів
    достатнього обсягу після сегментації);
    \item тематика --- міжнародні відносини, дипломатичні переговори,
    безпекова політика (що забезпечує однорідність дискурсивного стилю
    в межах корпусу);
    \item відсутність значних технічних перешкод (накладення голосів,
    музичний супровід, перебої зі звуком).
\end{enumerate}

Загалом для дослідження було відібрано 7 аудіозаписів щоденних звернень, датованих січнем---лютим
2025 року. Тематика записів охоплює дипломатичні зустрічі (тристоронні
зустрічі в ОАЕ), оборонну політику (доповіді Міністерства оборони, СБУ, ГУР),
енергетичну кризу (ліквідація наслідків ракетних ударів), міжнародну співпрацю
(переговори з лідерами США, Франції, Великої Британії, Норвегії, Австрії) та
внутрішню ситуацію (робота ремонтних бригад, соціальна підтримка).

\subsection{Підготовка корпусу}
\label{subsec:corpus_prep}

Підготовка корпусу складалася з п'яти послідовних етапів.

\subsubsection{Завантаження та конвертація аудіо}

Аудіозаписи було завантажено з офіційного каналу Офісу Президента за допомогою утиліти \texttt{yt-dlp} у форматі WebM (кодек Opus).
Оскільки формат WebM не підтримується безпосередньо бібліотеками обробки
звуку, кожен файл було конвертовано у формат WAV (PCM 16-bit, 44.1~kHz, моно)
за допомогою FFmpeg.

\subsubsection{Сегментація}

Для забезпечення оптимальної тривалості вхідних фрагментів для ASR-моделі
аудіофайли було розділено на сегменти тривалістю 3--5 хвилин за допомогою
спеціально розробленого скрипту на основі бібліотеки \texttt{pydub} (Python).
Алгоритм сегментації працює таким чином:
\begin{enumerate}[nosep]
    \item обчислюється RMS-енергія у ковзних вікнах;
    \item ідентифікуються паузи у мовленні --- ділянки, де рівень сигналу
    нижчий за поріг $-40$~дБFS протягом щонайменше 500~мс;
    \item обираються точки розрізки --- середини пауз, що забезпечують
    тривалість фрагментів у межах 3--5 хвилин;
    \item якщо залишковий фрагмент коротший за 1 хвилину, його приєднано
    до попереднього сегменту.
\end{enumerate}

Такий підхід гарантує, що розрізка відбувається на природних паузах мовлення,
а не посеред слова чи фрази. Загалом було отримано 12 аудіосегментів.
Перелік вихідних записів та результати сегментації наведено у
таблиці~\ref{tab:corpus}.

\begin{table}[h]
\centering
\caption{Склад корпусу дослідження}
\label{tab:corpus}
\small
\begin{tabular}{|c|l|c|c|}
\hline
\textnumero & Назва запису & К-сть сегм. & Прим. трив.\\
\hline
1 & В Еміратах сьогодні & 2 & $\approx$8 хв\\
2 & Діалог з Америкою & 2 & $\approx$8 хв\\
3 & Знаємо, що росіяни & 2 & $\approx$7 хв\\
4 & Командна ланка & 1 & $\approx$4 хв\\
5 & Наша стратегія & 2 & $\approx$7 хв\\
6 & Отримуємо хороші & 1 & $\approx$4 хв\\
7 & Поведінка & 2 & $\approx$7 хв\\
\hline
  & \textbf{Разом} & \textbf{12} & $\approx$\textbf{45 хв}\\
\hline
\end{tabular}
\end{table}

\subsubsection{Транскрипція}

Для автоматичної транскрипції використано модель Faster-Whisper medium
\cite{whisper} --- оптимізовану реалізацію моделі OpenAI Whisper на базі
CTranslate2, що забезпечує вищу швидкість інференсу при збереженні якості
розпізнавання. Параметри транскрипції наведено у таблиці~\ref{tab:whisper}.

\begin{table}[h]
\centering
\caption{Параметри транскрипції Faster-Whisper}
\label{tab:whisper}
\small
\begin{tabular}{|l|l|}
\hline
Параметр & Значення\\
\hline
Модель & medium\\
Мова & українська (\texttt{uk})\\
Beam size & 5\\
Температура & 0 (детерміністичний режим)\\
VAD-фільтр & увімкнено\\
Condition on previous text & увімкнено\\
\hline
\end{tabular}
\end{table}

Параметр \texttt{condition\_on\_previous\_text=True} було увімкнено для
забезпечення когерентності транскрипції в межах одного сегменту: модель
враховує попередньо згенерований текст при декодуванні наступних фрагментів,
що зменшує кількість контекстуальних помилок.

\subsubsection{Ручна корекція}

Після автоматичної транскрипції кожен текстовий файл було вручну перевірено
та відкориговано. Типовими помилками ASR-моделі виявились:
\begin{itemize}[nosep]
    \item спотворення власних назв та прізвищ (наприклад, <<Рустему Міров>>
    замість <<Рустем Умєров>>, <<Бодановим>> замість <<Будановим>>,
    <<Менуель>> замість <<Емманюель>>);
    \item помилки у географічних назвах (<<Валтавська область>> замість
    <<Полтавська область>>, <<Кривирі>> замість <<Кривий Ріг>>);
    \item лексичні підміни (<<вбойнів>> замість <<воїнів>>, <<упалення>>
    замість <<опалення>>);
    \item граматичні неточності.
\end{itemize}

Загалом у 12 сегментах було виявлено та виправлено близько 60--80 помилок
різного типу. Наявність цих помилок підтверджує доцільність ручної верифікації
транскрипцій перед етапом оцінювання якості машинного перекладу, оскільки
некоректна транскрипція каскадно впливає на всі подальші етапи конвеєра.

\subsubsection{Створення референсних перекладів}

Для кожного відкоригованого українського тексту було створено референсний
переклад англійською мовою, який виконує роль еталону (gold standard) для
порівняння з результатами машинного перекладу. Переклади було виконано за
допомогою великої мовної моделі Claude (Anthropic) із подальшою ручною
верифікацією, що включала перевірку фактичної коректності, точності передачі
власних назв та відповідності термінології. При перекладі дотримувались
таких вимог:
\begin{itemize}[nosep]
    \item збереження політико-дипломатичного регістру оригіналу;
    \item точність передачі власних назв, інституційних найменувань та
    військово-технічної термінології;
    \item збереження структури та порядку висловлювань для забезпечення
    коректного посегментного порівняння.
\end{itemize}

\subsection{Формат та структура даних}

Корпус організовано у дві паралельні папки: \texttt{Translations\_UA}
(українські тексти) та \texttt{Translations\_EN} (англійські референсні
переклади). Кожна папка містить 17 текстових файлів у кодуванні UTF-8:
\begin{itemize}[nosep]
    \item 12 сегментних файлів (за шаблоном
    \texttt{\{назва\}\_seg\{NNN\}.txt});
    \item 5 об'єднаних файлів (\texttt{\{назва\}\_full.txt}), що являють
    собою конкатенацію відповідних сегментів для записів, що складаються
    з двох або більше фрагментів.
\end{itemize}

Файли названо за першими словами кожного звернення, що забезпечує однозначну
відповідність між UA та EN версіями. Додатково у папці \texttt{segments/}
зберігаються відповідні аудіосегменти у форматі WAV.

Ця структура дозволяє на наступних етапах дослідження виконати машинний
переклад українських транскрипцій та порівняти результат із референсними
англійськими перекладами за обраними метриками.

\subsection{Висновки до розділу 3}

У третьому розділі описано процес формування корпусу дослідження. Корпус
складається з 12 аудіосегментів загальною тривалістю приблизно 45 хвилин,
отриманих із 7 звернень. Для кожного сегменту
підготовлено відкориговану транскрипцію українською мовою та референсний
переклад англійською. Описано алгоритм сегментації на основі пауз у
мовленні, параметри ASR-моделі Faster-Whisper medium, процедуру ручної
корекції транскрипцій та методику створення референсних перекладів.

\clearpage
